% Auto-split to keep each TeX source < 600 lines.
%
% (User input window)
%

\begin{proposition}[四阶含参核的唯一实谱分歧点：判别式五次因子与谱拓扑相变]\label{prop:pom-a4t-real-ep-branchpoint}
\leanverified{paper\_pom\_a4t\_real\_ep\_branchpoint}
沿注 \ref{rem:pom-collision-rh-break-a4-threshold} 的记号，令
\[
p_t(\lambda):=\det(\lambda I-A_4(t))=\lambda^{5}-2\lambda^{4}-t\lambda^{3}-2\lambda+2,
\qquad t\in\RR.
\]
则关于变量 \(\lambda\) 的判别式满足
\[
\mathrm{Disc}_\lambda(p_t)=-16\,Q(t),
\qquad
Q(t):=27t^5+35t^4+434t^3+4394t^2+10832t+7403.
\]
并且五次多项式 \(Q(t)\) 在实轴上仅有一个零点 \(t_{\mathrm{ep}}\in\RR\)（式 \eqref{eq:collision_kernel_A4_unit_circle_branchpoints}）。
因此 \(t=t_{\mathrm{ep}}\) 是该一参数族在实轴上唯一的谱分歧点：当且仅当 \(t=t_{\mathrm{ep}}\) 时，\(p_t\) 在 \(\RR\) 上出现二重特征值。
更进一步，谱的实/复根型在此发生跃迁：
\[
t>t_{\mathrm{ep}}
\ \Longrightarrow\
p_t\ \text{恰有三个实根与一对共轭非实根},
\qquad
t<t_{\mathrm{ep}}
\ \Longrightarrow\
p_t\ \text{恰有一个实根与两对共轭非实根}.
\]
\end{proposition}
\begin{proof}
判别式闭式已在注 \ref{rem:pom-a4t-discriminant-root-type} 给出；此处仅记其非平凡因子为 \(Q(t)\)。
Sturm 判别法给出 \(Q(t)\) 在 \(\RR\) 上恰有一个零点 \(t_{\mathrm{ep}}\)，从而实轴上的谱分歧点唯一（可复现脚本：\texttt{scripts/exp\_collision\_kernel\_A4\_unit\_circle\_branchpoints.py}）。
\par
当 \(t>t_{\mathrm{ep}}\) 时，\(Q(t)>0\)，故 \(\mathrm{Disc}_\lambda(p_t)<0\) 且 \(p_t\) 无重根。
对五次实系数多项式，判别式符号刻画非实共轭对数的奇偶性；由 \(\mathrm{Disc}_\lambda(p_t)<0\) 知其非实共轭对数为 \(1\)，因而根型为 \(3\) 实 \(+\) \(1\) 对非实共轭根。
\par
当 \(t<t_{\mathrm{ep}}\) 时，\(\mathrm{Disc}_\lambda(p_t)>0\) 且 \(p_t\) 无重根，故非实共轭对数为偶数。
取一点 \(t=-2<t_{\mathrm{ep}}\)，由 Sturm 计数可核验 \(p_{-2}\) 仅有一个实根；根型在连通区间 \((-\infty,t_{\mathrm{ep}})\) 上保持不变，遂得到 \(1\) 实 \(+\) \(2\) 对非实共轭根。证毕。
\end{proof}

% Auto-generated by scripts/exp_collision_kernel_A4_unit_circle_branchpoints.py
\begin{equation}\label{eq:collision_kernel_A4_unit_circle_branchpoints}
\begin{aligned}
Q(t)&:=27t^5+35t^4+434t^3+4394t^2+10832t+7403,\\
t_{\mathrm{ep}}&\approx -1.1943099264652641,\\
u_\ast(t)
&:=\frac{27t^{7}+62t^{6}+226t^{5}+1252t^{4}+3080t^{3}+5790t^{2}+7960t+4728}{27t^{8}+8t^{7}+264t^{6}+3056t^{5}+4160t^{4}-8064t^{3}-17728t^{2}-6832t+1984},\\
u_\ast(-1)&=1,\qquad u_\ast(1)=-1,\qquad u_\ast(t_{\mathrm{ep}})=:u_{\mathrm{ep}}\approx 0.6585687695847667,\\
\lambda_{\mathrm{ep}}&:=u_{\mathrm{ep}}^{-1/2}\approx 1.2322517229054710.
\end{aligned}
\end{equation}


\begin{proposition}[单位圆谱点的完全分类：仅三处实参数触碰 \texorpdfstring{$|\lambda|=1$}{|lambda|=1}]\label{prop:pom-a4t-unit-circle-spectrum-classification}
对实参数 \(t\)，方程 \(p_t(\lambda)=0\) 存在满足 \(|\lambda|=1\) 的根，当且仅当
\[
t\in\Bigl\{-\frac{11}{4},-1,1\Bigr\}.
\]
并且对应的单位圆根完全显式：
\[
t=-1\ \Longleftrightarrow\ \lambda\in\{1,-1\},
\qquad
t=1\ \Longleftrightarrow\ \lambda\in\{\mathrm{i},-\mathrm{i}\},
\]
\[
t=-\frac{11}{4}
\ \Longleftrightarrow\
\lambda\in\left\{\frac{3\pm \mathrm{i}\sqrt7}{4}\right\}.
\]
其中 \(t=-\frac{11}{4}\) 给出唯一的“非割圆”单位圆触点：单位圆根生成二次虚域 \(\QQ(\sqrt{-7})\)，并以共轭对出现在谱中。
\end{proposition}
\begin{proof}
设 \(t\in\RR\) 且 \(|\lambda|=1\) 满足 \(p_t(\lambda)=0\)。
由于 \(p_t\) 系数全实，有 \(0=\overline{p_t(\lambda)}=p_t(\overline\lambda)\)；又因 \(|\lambda|=1\)，\(\overline\lambda=\lambda^{-1}\)，故亦有 \(p_t(1/\lambda)=0\)。
\par
将 \(p_t(\lambda)=0\) 与 \(\lambda^5p_t(1/\lambda)=0\) 视为关于 \(t\) 的两条方程消元，得到
\[
(\lambda^4-1)\,(2\lambda^2-3\lambda+2)=0.
\]
因此 \(|\lambda|=1\) 的可能值仅为 \(\lambda\in\{\pm1,\pm\mathrm{i}\}\) 或满足 \(2\lambda^2-3\lambda+2=0\) 的一对共轭根 \(\lambda=(3\pm \mathrm{i}\sqrt7)/4\)。
将其代回
\(
t=\frac{\lambda^5-2\lambda^4-2\lambda+2}{\lambda^3}
\)
即可分别得到 \(t=-1,1,-11/4\)。
反向蕴含由直接代入验证。证毕。
\end{proof}

\leanverified{paper\_pom\_a4t\_unit\_circle\_spectrum\_classification}

\begin{corollary}[偶谱的单位圆触点：\texorpdfstring{$\det(I-uA_4(t)^2)$}{det(I-uA4(t)^2)} 的三类单位模零点]\label{cor:pom-a4t-even-spectrum-unit-circle-touchpoints}
令 \(E_t(u):=\det(I-uA_4(t)^2)\) 如命题 \ref{prop:pom-a4t-even-zeta-quintic}。
则对 \(t\in\RR\)，方程 \(E_t(u)=0\) 在单位圆 \(|u|=1\) 上出现零点，当且仅当
\[
t\in\Bigl\{-\frac{11}{4},-1,1\Bigr\}.
\]
并且对应零点可取为
\[
t=-1\Rightarrow u=1,
\qquad
t=1\Rightarrow u=-1,
\qquad
t=-\frac{11}{4}\Rightarrow 4u^{2}-u+4=0,
\]
即 \(u=(1\pm 3\mathrm{i}\sqrt7)/8\in\QQ(\sqrt{-7})\) 给出一对单位模偶谱零点。
\end{corollary}
\begin{proof}
由谱映射定理，
\(
E_t(u)=0
\)
当且仅当存在 \(\lambda\in\mathrm{spec}(A_4(t))\) 使得 \(u=\lambda^{-2}\)。
因此 \(|u|=1\) 当且仅当存在 \(|\lambda|=1\) 的特征值。
结合命题 \ref{prop:pom-a4t-unit-circle-spectrum-classification} 得到参数集合为 \(\{-11/4,-1,1\}\)。
三类 \(u\) 的显式值由 \(u=\lambda^{-2}\) 直接计算，并对 \(t=-11/4\) 的情形化为二次式 \(4u^2-u+4=0\)。证毕。
\end{proof}

\begin{proposition}[偶谱二重根的统一有理表达：三处退化共享同一线性子结果式]\label{prop:pom-a4t-even-spectrum-double-root-ustar}
令 \(E_t(u)\) 如命题 \ref{prop:pom-a4t-even-zeta-quintic}，并令 \(t_{\mathrm{ep}}\) 为命题 \ref{prop:pom-a4t-real-ep-branchpoint} 的唯一实分歧点。
则存在显式有理函数 \(u_\ast(t)\)（式 \eqref{eq:collision_kernel_A4_unit_circle_branchpoints}），满足下列等价刻画：
\[
\text{对 }t\in\RR,\qquad
\Bigl(\exists\,u\in\CC:\ E_t(u)=0\ \text{且}\ \partial_uE_t(u)=0\Bigr)
\quad\Longleftrightarrow\quad
t\in\{-1,1,t_{\mathrm{ep}}\},
\]
且在这些退化参数上，\(u=u_\ast(t)\) 为 \(E_t(u)\) 的二重根，并满足
\[
u_\ast(-1)=1,\qquad u_\ast(1)=-1,\qquad u_\ast(t_{\mathrm{ep}})=:u_{\mathrm{ep}}\approx 0.6585687696.
\]
因此在唯一实谱分歧点 \(t=t_{\mathrm{ep}}\) 处，\(A_4(t)\) 出现二重特征值 \(\lambda_{\mathrm{ep}}\in\RR_{>0}\) 使得
\(
u_{\mathrm{ep}}=\lambda_{\mathrm{ep}}^{-2}
\)
（数值上 \(\lambda_{\mathrm{ep}}\approx 1.2322517229\)，见式 \eqref{eq:collision_kernel_A4_unit_circle_branchpoints}）。
\end{proposition}
\begin{proof}
由式 \eqref{eq:collision_kernel_A4_even_zeta_quintic}，
\(\mathrm{Disc}_u(E_t)=-256(t-1)^2(t+1)^2Q(t)\)。
因此 \(E_t\) 在 \(u\)-变量上出现重根当且仅当 \(t=\pm1\) 或 \(Q(t)=0\)。
而命题 \ref{prop:pom-a4t-real-ep-branchpoint} 表明 \(Q(t)=0\) 在实轴上仅对应 \(t=t_{\mathrm{ep}}\)，从而得到退化参数集合 \(\{-1,1,t_{\mathrm{ep}}\}\)。
\par
在 \(\ZZ[t]\) 系数域上对 \((E_t,\partial_uE_t)\) 作 Euclid/子结式消元，其次数为 \(1\) 的子结式给出一条线性因子
\(
(\text{多项式}(t))\cdot u-(\text{多项式}(t))
\)；
对应解即为式 \eqref{eq:collision_kernel_A4_unit_circle_branchpoints} 中的 \(u_\ast(t)\)，并满足
\(
E_t(u_\ast(t))=\partial_uE_t(u_\ast(t))=0
\)
在上述退化参数处同时成立（可复现脚本同命题 \ref{prop:pom-a4t-real-ep-branchpoint}）。
\par
最后，对 \(t=t_{\mathrm{ep}}\) 的情形，\(\lambda\)-层的二重根由 \(p_t\) 的判别式零化给出；
由 \(E_t(u)=\prod_{\lambda\in\mathrm{spec}(A_4(t))}(1-u\lambda^2)\) 得到 \(u_{\mathrm{ep}}=\lambda_{\mathrm{ep}}^{-2}\)。
数值近似由式 \eqref{eq:collision_kernel_A4_unit_circle_branchpoints} 给出。证毕。
\end{proof}

\leanverified{paper\_pom\_a4t\_even\_spectrum\_unit\_circle\_touchpoints}

\leanverified{paper\_pom\_a4t\_even\_spectrum\_double\_root\_ustar}

\begin{remark}[Newman 临界与单位圆几何严格分离：RH 越线发生在单位圆内相]\label{rem:pom-a4t-newman-unit-circle-separation}
由命题 \ref{prop:pom-a4t-real-ep-branchpoint} 与式 \eqref{eq:collision_kernel_A4_unit_circle_branchpoints}，有
\(
-\frac{11}{4}<t_{\mathrm{ep}}<-1<1
\)；
另一方面，注 \ref{rem:pom-collision-rh-break-a4-threshold} 给出四阶 Newman 临界阈值 \(t_\star\) 为唯一正实解，且 \(t_\star\approx 6.9757>1\)（式 \eqref{eq:collision_kernel_A4_edge_weight_threshold}）。
因此严格次序链成立：
\[
\boxed{\;-\frac{11}{4}<t_{\mathrm{ep}}<-1<1<t_\star\;}
\]
从而 \(t_\star\) 处的核版 RH 越线发生在全谱根型已固定的解析支内部（\(3\) 实 \(+\) \(1\) 对非实共轭），并且位于单位圆触点 \(t=1\) 的右侧。
换言之，临界越线可在纯实谱通道（Perron 根与主负实根）的模长竞争框架内刻画，而与 \(|\lambda|=1\) 触碰及谱分歧机制在几何上严格正交。
\end{remark}

\begin{conclusion}[四阶核族的分歧分层：全谱唯一实异常点与偶谱三处二重退化]\label{con:pom-a4t-branching-layered-structure}
对实参数 \(t\)，全谱多项式 \(p_t(\lambda)=\det(\lambda I-A_4(t))\) 的实轴谱分歧仅发生在 \(t=t_{\mathrm{ep}}\)（命题 \ref{prop:pom-a4t-real-ep-branchpoint}），因此除该点外五条特征值分支在 \(t\) 上均为实解析函数。
而偶谱多项式 \(E_t(u)=\det(I-uA_4(t)^2)\) 的判别式在 \(t=\pm1\) 与 \(t=t_{\mathrm{ep}}\) 处同时退化（式 \eqref{eq:collision_kernel_A4_even_zeta_quintic} 与命题 \ref{prop:pom-a4t-even-spectrum-double-root-ustar}）：
其中 \(t=\pm1\) 的二重根来自平方映射 \(\lambda\mapsto \lambda^2\) 的识别退化，
而 \(t=t_{\mathrm{ep}}\) 的二重根对应真实的谱分歧（\(\lambda\)-层二重特征值）。
此外，Newman 临界 \(t_\star\) 位于无分歧解析支内部（注 \ref{rem:pom-collision-rh-break-a4-threshold}），属于纯粹的模长竞争相变而非谱分歧相变。
\end{conclusion}
